%% From paper/abstract_arxiv.txt (1,876 chars, within the arXiv 1,920 cap).
Agentic workloads interleave LLM calls with tool execution, and when a session returns from a tool call sets its serving cost. On an RBLN CA25 NPU running vLLM, we decompose that return-arrival process on real hardware into three independent mechanisms: alignment between the concurrent request count and the compiled decode-batch grid, the reuse cliff of a sequence-granular FIFO slot pool, and the serialisation tax of exclusively executed prefill. None responds to a property of an individual request; all three respond to a collective property of arrivals. We establish the first causally, with a recompile intervention that changes only the grid. A step-level simulator carrying all three, with zero fitted parameters, passes a preregistered out-of-sample gate.

On that model we obtain a negative result. Rescheduling returns leaves real offline headroom, but a median 60\% of it is coordination: hand every session omniscient knowledge and let it decide independently, and that share disappears. No per-session runtime policy can reach it, in principle. Nor is the remainder reachable: the value of return-time information is workload-specific and vanishes under measured tool latencies, and a published duration estimator, borrowed verbatim, misses the accuracy threshold by a margin four more orders of sample magnitude do not close.

That leaves one reachable lever: a compile-time configuration, which is coordination decided once, in advance, for everybody. Escapement chooses one from workload distribution statistics and the uncalibrated simulator alone, with no device measurement in the selection loop, applies it with a seven-minute recompile, and confirms a 9.7-10.1\% device-time recovery at N=8 concurrent sessions, on two channels registered before measurement. Some levers can be bought with a distribution even when no one can predict the individual draw.
